\documentclass[11pt]{article}
\usepackage[margin=1.1in]{geometry}
\usepackage{amsmath,amssymb,amsthm}
\usepackage{booktabs}
\usepackage{url}
\usepackage[hidelinks]{hyperref}

\newtheorem{theorem}{Theorem}
\newtheorem{lemma}[theorem]{Lemma}
\newtheorem{proposition}[theorem]{Proposition}
\newtheorem{corollary}[theorem]{Corollary}
\theoremstyle{remark}
\newtheorem{remark}[theorem]{Remark}
\theoremstyle{definition}
\newtheorem{openproblem}[theorem]{Open Problem}
\newtheorem{conjecture}[theorem]{Conjecture}

\newcommand{\frc}[1]{\{#1\}}
\newcommand{\fl}[1]{\lfloor #1\rfloor}
\newcommand{\M}{\mathsf{M}}
\newcommand{\Ot}{\widetilde{O}}
\newcommand{\Z}{\mathbb{Z}}

\title{Extracting decimal digits of $\pi$ in quasi-linear time\\ from binary BBP formulas}
\author{Joseph Bab\v{c}anec\\ \small Benedict College\\ \small \texttt{joseph.babcanec@benedict.edu}}
\date{July 2026}

% NOTE (not compiled): consider adding an AI-assistance disclosure in the
% acknowledgments if the target venue requires one.

\begin{document}
\maketitle

\begin{abstract}
The Bailey--Borwein--Plouffe (BBP) formula computes isolated binary or
hexadecimal digits of $\pi$ without computing earlier digits, but no
analogous decimal formula exists: Borwein, Borwein and Galway proved that
$\pi$ admits no Machin-type BBP arctangent formula in any base that is not
a power of $2$, and every published decimal digit-extraction algorithm
runs in quadratic time or worse. We show that the $N$-th decimal digit of
$\pi$ can nevertheless be computed from the ordinary hexadecimal BBP
formula in quasi-linear time --- $O(N\log^{3}N)$ bit operations --- at the
price of $\Theta(N)$ bits of read-only auxiliary space holding the binary
expansion of the single integer $5^{\,N-1}$: after multiplying the BBP
series termwise by $10^{N-1}$, every term takes the form
$\pm2^{E}5^{\,N-1}/\mu$ with $\mu$ odd; terms with $E\ge0$ reduce to
word-size modular exponentiation, while the $\Theta(N)$ terms with $E<0$
satisfy an exact identity that expresses their fractional parts through
residues and bit-windows of the \emph{shared} integer $5^{\,N-1}$,
batchable with accumulating remainder trees. The method applies verbatim
to every constant with a BBP formula in a base $2^{r}$ (e.g.\ $\ln2$),
yields the digits of $\pi$ in \emph{any} base $B$ in $\Ot(N)$ time, and
resolves the obstruction described by Zudilin (arXiv:2409.10097), whose
base-5 (originally base-10) attempt fails precisely on the $E<0$ terms.
Implementations in Python and C++ are verified at decimal positions
through $10^{7}$ against independent computations and public digit
corpora. We then reduce every phase but one to $O(N\log\log N)$ word
operations: sieving the modulus progressions turns all modular
exponentiations into geometric progressions along prime-indexed
subprogressions, with $O(1)$ starters by Euler's criterion --- an
amortization that applies verbatim to classical binary BBP extraction ---
while a B\'ezout split makes half of every remaining term word arithmetic
and the other halves consolidate into a \emph{single} $2$-adic integer
$V$ with a division-free construction. A GMP-backed implementation of
this refined pipeline measures a fitted exponent of $N^{1.08}$ across
positions $10^{4}$--$10^{7}$ --- the quasi-linear regime in practice,
thirty times faster at $10^{7}$ than a multithreaded implementation of
the unbatched algorithm. The one remaining logarithm is analyzed as a
barrier with seven machine-validated faces --- merge depth,
rising-factorial products, the discrete-logarithm sum
$E=\sum_{k\le K}\operatorname{dlog}_3(1+8k)$, a $2$-adically
\emph{convergent} Euler--Maclaurin identity linking $E$ to Diamond's
$2$-adic log-gamma function, and an elementary huge-exponent power sum
--- and a classical factorization technique of Borwein then splits them:
the products, discrete-logarithm sums, power sums, and log-gamma values
all fall to $O(N\log^{2}N)$ (machine-validated), while the digit
pipeline's additive core --- one consolidated sum of scaled $2$-adic
inverses --- remains at the product-tree bound and is isolated as the
precise open problem. No factoring-hardness argument protects it. We
state the conjectural $N\log N$ floor precisely, together with the
sublinear-time and polylogarithmic-space questions that remain open.
\end{abstract}

\section{Introduction}

Write $d_N(\xi)$ for the $N$-th decimal digit of $\xi>0$ after the decimal
point, so that with $\frc{x}=x-\fl{x}$,
\begin{equation}\label{eq:digit}
d_N(\xi) \;=\; \bigl\lfloor 10\,\frc{10^{\,N-1}\xi}\bigr\rfloor .
\end{equation}
The \emph{digit extraction problem} asks for $d_N(\pi)$ without computing the
$N-1$ digits before it. Bailey, Borwein and Plouffe~\cite{BBP97} solved the
binary case with the formula
\begin{equation}\label{eq:bbp}
\pi \;=\; \sum_{k=0}^{\infty}\frac{1}{16^{k}}
\Bigl(\frac{4}{8k+1}-\frac{2}{8k+4}-\frac{1}{8k+5}-\frac{1}{8k+6}\Bigr),
\end{equation}
which yields hexadecimal digits at position $N$ in $O(N\log N)$-type time and
$O(\log N)$ space via per-term modular exponentiation. For base $10$ no such
formula exists: Borwein, Borwein and Galway~\cite{BBG04} proved that $\pi$ has
no Machin-type BBP arctangent formula in any base $b\neq 2^{m}$. The known
decimal algorithms form a ladder of quadratic-or-worse methods:
Plouffe's original scheme~\cite{Plouffe96} runs in $O(N^{3}\log^{3}N)$ time,
Bellard's improvement~\cite{Bellard97} in $O(N^{2})$ time in any base, and
Gourdon's method~\cite{Gourdon03} in $O\!\bigl(N^{2}\log\log N/\log^{2}N\bigr)$
word operations with $O(\log^{2}N)$ memory. Plouffe's 2022
approach~\cite{Plouffe22} extracts a digit from an asymptotic expansion whose
coefficients are Bernoulli (or Euler) numbers; it requires precomputed tables
of those numbers and is thereby capped near position $10^{8}$, the current
extent of Bernoulli computations. Most recently
Zudilin~\cite{Zudilin24} attempted a BBP-style computation of $\pi$ in base
$5$ --- the first arXiv version was titled ``\ldots in base 10'' --- via a
Salikhov-type representation over $\mathbb{Z}[i]$, and identified an
``obvious flaw'': his per-term modular identity fails exactly on the terms in
which the power of $5$ in the numerator is exhausted, and no truncation
avoids them. He left the repair as an open question; the note has no
published successor.

This paper repairs that flaw --- not for Zudilin's Gaussian series
specifically, but for the ordinary formula~\eqref{eq:bbp} and, with it, every
binary BBP constant. The price is memory: we abandon the $O(\log N)$-space
regime of BBP and store the binary expansion of the single integer
$X=5^{\,N-1}$, about $0.2902\,N$ bytes. In exchange the running time becomes
quasi-linear.

\begin{theorem}[main]\label{thm:main}
Let $\M(n)$ denote the cost of multiplying two $n$-bit integers. The $N$-th
decimal digit of $\pi$ (together with any constant window of following
digits) can be computed, without computing any earlier digit, in
\[
O\bigl(\M(N) \;+\; \M(N\log N)\log N \;+\; N\log^{2}N\log\log N\bigr)
\;\subseteq\; O\bigl(N\log^{3}N\bigr)
\]
bit operations, using $\Theta(N)$ bits of auxiliary space of which all but
$O(\log^{2} N)$ bits are a read-only string shared by all terms, namely the
binary expansion of $5^{\,N-1}$. The output is certified in the Las Vegas
sense of Remark~\ref{rem:guard}.
\end{theorem}

The same statement holds with $\pi$ replaced by any constant possessing a
BBP-type formula in a base $2^{r}$ (Theorem~\ref{thm:general}), for instance
$\ln 2$; and with ``decimal'' replaced by ``base $B$'' for arbitrary $B\ge 2$
(Corollary~\ref{cor:anybase}), which supersedes the $O(N^{2})$ any-base
algorithm of~\cite{Bellard97} in time at the stated memory cost.
Section~\ref{sec:improve} then sharpens the algorithm structurally: all
phases except one collapse to $O(N\log\log N)$ word operations, and the
surviving phase --- the construction of a single $2$-adic integer --- is
analyzed as a barrier in Section~\ref{sec:barrier}.

\paragraph{What is, and is not, being claimed.}
Computing \emph{all} $N$ digits of $\pi$ by binary splitting of the
Chudnovsky series already takes softly linear time, so quasi-linear cost per
isolated digit is not by itself a new capability; and the canonical open
problem posed in~\cite{BBP97} --- decimal digit extraction in polylogarithmic
(or even sublinear) space --- remains open and is \emph{not} resolved here,
since our space is $\Theta(N)$ bits. The contributions are:
(i) the first digit-extraction algorithm for decimal $\pi$ based on a BBP
formula, with per-term independence: after the one shared precomputation the
terms are processed in any order by any number of workers that communicate
only word-size results, in the style of distributed binary-BBP projects;
(ii) to our knowledge the first \emph{implemented and verified} subquadratic
decimal extractor in any model (Section~\ref{sec:impl});
(iii) the resolution of the specific obstruction of~\cite{Zudilin24};
(iv) an auxiliary object --- one power of $5$ computed on the fly --- rather
than tables that must be precomputed by other means, in contrast
to~\cite{Plouffe22};
(v) the sieve-batched exponentiation of Section~\ref{sec:improve}, which
also accelerates classical binary BBP extraction; and
(vi) a precise, machine-validated characterization of the obstruction to
going faster (Section~\ref{sec:barrier}).
Closest to (ii) at the level of results, Gourdon's unpublished 2003
manuscript~\cite{Gourdon03} asserts in its final section a time--memory
tradeoff, $O\!\bigl(n^{2}\log^{3}n\log\log n/(m\log^{2}(n/m))\bigr)$ time at
memory $m\in[n^{\varepsilon},n]$, whose endpoint $m\approx n$ would likewise
be softly linear; the manuscript states that ``the details of the algorithm
will be added soon in a next version,'' but no later version ever appeared
(the archived file is unchanged from February 2003 through 2026), no details
or implementation are known, and the sketch that is given (binary splitting
over blocks) differs structurally from the present method. We therefore claim
the algorithm and its analysis, not priority over the bare assertion that
some such tradeoff should exist.

\section{The three-phase decomposition}\label{sec:method}

Throughout, $N\ge 1$ is the target position, $n:=N-1$, $X:=5^{\,n}$, and
$L:=\lceil n\log_2 5\rceil$ the bit length of $X$. Since
$x\mapsto\frc{x}$ is a homomorphism $\mathbb{R}\to\mathbb{R}/\mathbb{Z}$,
fractional parts of individual terms may be accumulated with signs modulo
$1$; alternating or negative coefficients pose no obstacle.

\subsection{Normal form}

Multiply \eqref{eq:bbp} termwise by $10^{\,n}=2^{\,n}5^{\,n}$ and absorb all
powers of two --- the numerators $4,2,1,1$ and the even parts of
$8k+4=4(2k+1)$ and $8k+6=2(4k+3)$ --- into a single exponent. With
$A:=n-4k$:

\begin{lemma}[normal form]\label{lem:normal}
\begin{equation}\label{eq:normal}
10^{\,n}\pi \;=\; \sum_{k\ge 0}\Bigl[
\;2^{A+2}\frac{X}{8k+1}
\;-\;2^{A-1}\frac{X}{2k+1}
\;-\;2^{A}\frac{X}{8k+5}
\;-\;2^{A-1}\frac{X}{4k+3}\Bigr],
\end{equation}
i.e.\ every term equals $\pm\,2^{E} X/\mu$ with $\mu$ odd and
$E\in\{A+2,\,A-1,\,A,\,A-1\}$.
\end{lemma}

\begin{proof}
Immediate from $10^{\,n}16^{-k}=2^{\,n-4k}5^{\,n}$ and the factorizations
$8k+4=2^{2}(2k+1)$, $8k+6=2(4k+3)$.
\end{proof}

Terms decay like $16^{-k}$, so positions $N$ through $N+O(1)$ are determined
by the terms with $k\le K$, where
$K=\lceil (n\log_2 10 + G + c)/4\rceil$ for a $G$-bit accumulator
(Section~\ref{sec:error}); in total $4(K+1)\approx 3.33\,N$ terms. Each term
falls into one of three regimes.

\subsection{Phase 1: easy terms ($E\ge 0$)}

Here $2^{E}X$ is a nonnegative integer, so
\[
\Bigl\{\frac{2^{E}X}{\mu}\Bigr\}
=\frac{\bigl(2^{E}\bmod \mu\bigr)\bigl(5^{\,n}\bmod \mu\bigr)\bmod \mu}{\mu},
\]
computable with two modular exponentiations modulo the word-size integer
$\mu$ --- exactly as in classical BBP. There are $\Theta(N)$ such terms
($E\ge0$ holds for $k\lesssim n/4$ in each family).

\subsection{Phase 2: band terms ($E<0$, $J:=-E< L$)}

These are the terms on which every previous decimal approach loses: the
numerator is no longer an integer, the natural modulus $\mu 2^{J}$ grows
exponentially along the family, and per-term reduction of $X$ modulo
$\mu2^J$ costs $\Theta(\M(N))$, giving $\Theta(N)$-many big operations ---
the quadratic wall, and precisely the failure identified
in~\cite[\S3]{Zudilin24}. The wall disappears because \emph{all band terms
share the same numerator $X$}:

\begin{lemma}[band identity]\label{lem:band}
Let $X,\mu\ge 1$ and $J\ge 1$ be integers, $H:=\fl{X/2^{J}}$,
$h:=H\bmod \mu$, and $u:=\frc{X/2^{J}}\in[0,1)$. Then
\[
\Bigl\{\frac{X}{\mu\,2^{J}}\Bigr\} \;=\; \frac{h+u}{\mu}.
\]
\end{lemma}

\begin{proof}
$X/(\mu2^{J})=(H+u)/\mu$. Writing $H=q\mu+h$,
$X/(\mu 2^{J})=q+(h+u)/\mu$ with $0\le (h+u)/\mu<(\mu-1+1)/\mu=1$, so
$(h+u)/\mu$ is the fractional part.
\end{proof}

\begin{lemma}[recovering $h$]\label{lem:h}
Let $\mu$ be odd, $R:=X\bmod\mu$ and
$\widetilde{Y}:=\bigl(X\bmod 2^{J}\bigr)\bmod\mu$. Then
\[
h \;=\; \bigl(R-\widetilde{Y}\bigr)\cdot 2^{-J} \;\bmod\; \mu .
\]
\end{lemma}

\begin{proof}
From $X=2^{J}H+(X\bmod 2^{J})$, reduce modulo $\mu$ and use that $2$ is
invertible mod $\mu$.
\end{proof}

Consequently a band term needs:
$R=5^{\,n}\bmod\mu$ and $2^{-J}\bmod\mu=((\mu+1)/2)^{J}\bmod\mu$, two
word-size modular exponentiations; the \emph{prefix residue}
$\widetilde{Y}$, discussed in Section~\ref{sec:batch}; and $u$ to $G_p$ bits,
which is by definition the bit-window of $X$ at depth $J$:
\[
\bigl\lfloor 2^{G_p}u\bigr\rfloor
=\Bigl\lfloor \frac{X\bmod 2^{J}}{2^{\,J-G_p}}\Bigr\rfloor
=\text{bits } [\,J-G_p,\;J\,) \text{ of } X,
\]
an $O(1)$ read from the stored binary expansion of $X$. This is the entire
role of the $\Theta(N)$-bit auxiliary string: it converts every band term's
``impossible'' modulus $\mu2^{J}$ into one shared object read locally.

\subsection{Phase 3: tail terms ($J\ge L$)}

Then $H=0$, so the term is $u/\mu$ with
$u=X/2^{J}<1$ read from the top of the stored string; only
$O(G)$ terms lie between $J=L$ and the accumulator cutoff.

\subsection{Accumulation and error}\label{sec:error}

All contributions are accumulated in a fixed-point register of $G$ fraction
bits, each term entering as
$\fl{2^{G}(h+u)/\mu}$ (band/tail) or $\fl{2^{G}r/\mu}$ (easy), with sign,
modulo $2^{G}$.

\begin{proposition}[error bound]\label{prop:error}
With window precision $G_p$ and $T$ terms processed, the accumulator differs
from $2^{G}\frc{10^{\,n}\pi}$ by at most
$T\bigl(1+2^{\,G-G_p}\bigr)+2$ units, i.e.\ by
$T\cdot 2^{-G}+T\cdot 2^{-G_p}+O(2^{-G})$ after rescaling. Choosing
$G_p,G=\Theta(\log N)$ with $G_p\ge \log_2 T+g$ and $G\ge G_p+48$ leaves
$g$ trustworthy guard bits above the output window.
\end{proposition}

\begin{proof}
Each floor loses less than one unit in the last place; truncating $u$ to
$G_p$ bits perturbs $(h+u)/\mu$ by less than $2^{-G_p}/\mu\le 2^{-G_p}$; the
series cutoff contributes less than one unit by the choice of $K$. Sum over
$T$ terms.
\end{proof}

\begin{remark}[certification]\label{rem:guard}
As with all BBP-type extraction, the digit window is read off the top of the
accumulator and is provably correct unless $\frc{10^{\,n}\pi}$ lies within
the total error $\varepsilon$ of a carry boundary of the window --- i.e.\
unless the digits following the window form a run of $9$s or $0$s of length
$\approx(\log_{10}\varepsilon^{-1})-w$. The algorithm detects this condition
in the output and re-runs with doubled $G$; under the standard heuristic that
the digits of $\pi$ are equidistributed the expected number of passes is
$1+o(1)$. (An unconditional guard follows from the irrationality measure
$\mu(\pi)\le 7.11$ of Zeilberger--Zudilin~\cite{ZZ15}, but only with
$G=\Theta(N)$, which would sacrifice quasi-linearity; every implementation
of binary BBP makes the same tradeoff.)
\end{remark}

\section{Complexity}\label{sec:complexity}

\subsection{Precomputation and per-term costs}

$X=5^{\,n}$ is computed by repeated squaring in $O(\M(N))$ bit operations
--- the $O(\log N)$ squarings act on geometrically growing operands, so the
cost telescopes --- and stored as $\approx 0.2902N$ bytes. The $O(N)$ modular
exponentiations (Phases 1--2) cost $O(\log N)$ multiplications of
$O(\log N)$-bit words each: $O(N\log N)$ word operations, i.e.\
$O(N\log^{2}N\log\log N)$ bit operations; window reads are $O(1)$ words
each.

\subsection{Batched prefix residues}\label{sec:batch}

The one nontrivial cost is the family of prefix residues
$\widetilde{Y}_t=(X\bmod 2^{J_t})\bmod\mu_t$ over the $T=\Theta(N)$ band
terms. Computed naively (a Horner pass over the stored string per term) each
costs $O(J_t/w)$ word operations --- an $O(N^{2}/w)$ total that our v1
implementations accept (Section~\ref{sec:impl}). It is however batchable
with remainder-tree technology~\cite{BorodinMoenck74,Bernstein04}:

\begin{theorem}[prefix-residue batching]\label{thm:batch}
Given the binary expansion of $X$ (of length $L$) and pairs
$(J_1,\mu_1),\dots,(J_T,\mu_T)$ with word-size moduli and the $J_t$ in
arithmetic progression, all residues $(X\bmod 2^{J_t})\bmod \mu_t$ can be
computed in $O\bigl(\M(L+Tw)\,\log T\bigr)$ bit operations.
\end{theorem}

\begin{proof}[Proof sketch]
Let $\delta$ be the common difference of the $J_t$ (for the four families
of~\eqref{eq:normal}, $\delta=4$; run one instance per family) and let
$c_t\in[0,2^{\delta})$ be the bits of $X$ in $[J_t,J_{t+1})$. The prefix
values $P_t:=X\bmod 2^{J_t}$ obey the linear recurrence
$P_{t+1}=P_t+2^{J_t}c_t$, i.e.\ the pair $v_t=(P_t,\,2^{J_t})^{\mathsf T}$
transforms as $v_{t+1}=A_t\,v_t$ with the word-size upper-triangular matrix
\[
A_t=\begin{pmatrix}1 & c_t\\ 0 & 2^{\delta}\end{pmatrix},
\qquad\text{so}\qquad
\widetilde{Y}_t=\bigl[(A_{t-1}\cdots A_0)\,v_0\bigr]_1 \bmod \mu_t .
\]
Computing a running product of word-size matrices reduced modulo a
different word-size modulus at each leaf is precisely the
\emph{accumulating remainder tree} of Costa--Gerbicz--Harvey and
Harvey--Sutherland \cite{CGH14,HS14}: one product/remainder tree over the
$T$ leaves in which the partial products (which grow to $\Theta(L)$ bits at
the root) are pushed down and reduced modulo the subtree modulus products
(total size $\Theta(Tw)$). Each of the $O(\log T)$ levels costs
$O(\M(L+Tw))$ by superadditivity of $\M$, giving the bound. The workspace is
$O(L+Tw)$ bits, processed level by level.
\end{proof}

With $L=\Theta(N)$, $T=\Theta(N)$, $w=O(\log N)$ this is
$O(\M(N\log N)\log N)\subseteq O(N\log^{3}N)$, and Theorem~\ref{thm:main}
follows by summing the phase costs: the accumulating tree dominates, the
modular exponentiations contribute $O(N\log^{2}N\log\log N)$ bits, and
everything else is smaller. In the word RAM the totals are
$O(N\log^{2}N)$ word operations for the tree and $O(N\log N)$ for the
exponentiations --- and Section~\ref{sec:improve} reduces the latter, and
every other non-tree cost, to $O(N\log\log N)$.

\subsection{Comparison}

\begin{center}
\footnotesize
\setlength{\tabcolsep}{4pt}
\begin{tabular}{lllll}
\toprule
Method & Time & Space & Base & Status\\
\midrule
BBP 1997~\cite{BBP97} & $\Ot(N)$ & $O(\log N)$ & $2^{r}$ only & implemented\\
Plouffe 1996~\cite{Plouffe96} & $O(N^{3}\log^{3}N)$ & $O(\log N)$ & any & implemented\\
Bellard 1997~\cite{Bellard97} & $O(N^{2})$ & $O(\log N)$ & any & implemented\\
Gourdon 2003~\cite{Gourdon03}, Thm.~1 & $O(N^{2}\log\log N/\log^{2}N)$\,$^{\mathrm a}$ & $O(\log^{2}N)$ & $10$ & implemented\\
Gourdon 2003~\cite{Gourdon03}, Thm.~2 & $O\!\bigl(\tfrac{N^{2}\log^{3}N\log\log N}{m\log^{2}(N/m)}\bigr)$ & $O(m)$ & $10$ & asserted only\\
Plouffe 2022~\cite{Plouffe22} & (no analysis; $N\le 10^{8}$) & tables $B_{2n}$ & $10$, $2$ & implemented\\
\textbf{This paper} & $O(N\log^{3}N)$ & $\Theta(N)$ bits (shared, r/o) & any & implemented$^{\mathrm b}$\\
\bottomrule
\end{tabular}

\smallskip
\parbox{0.92\textwidth}{\footnotesize $^{\mathrm a}$word operations.\quad
$^{\mathrm b}$measured $N^{1.077}$ over $N\in[10^{4},10^{7}]$ by the
refined pipeline (v3); see Section~\ref{sec:impl}.}
\end{center}

\section{Faster phases: sieve-batching and a $2$-adic normal form}
\label{sec:improve}

\subsection{Sieve-batched exponentiation}\label{sec:sieve}

The exponentiation workload of Sections \ref{sec:method}--\ref{sec:complexity}
consists, per term, of powers of $2$ whose exponents are \emph{linear in the
index} $k$ (namely $2^{\pm4k}$-type factors and $2^{-J}$ with $J$ linear in
$k$) and powers with \emph{fixed} exponents ($5^{\,n}$ and $2^{\,n+O(1)}$
modulo $\mu$). Square-and-multiply prices each at $O(\log N)$ word
multiplications: $O(N\log N)$ in total. Both kinds amortize.

\begin{lemma}[Euler starters]\label{lem:euler}
Let $q$ be an odd prime and $k_0$ minimal with $q\mid 8k_0+1$, and write
$8k_0+1=qs$; then $s\le 7$ is odd and
\[
2^{\,4k_0}\;\equiv\;\Bigl(\tfrac{2}{q}\Bigr)^{\!s}\,2^{(s-1)/2}
\pmod q ,
\]
computable in $O(1)$ word operations. Analogous $O(1)$ formulas hold for
the progressions $8k+5$, $4k+3$, $2k+1$ (using
$2^{\,qs-3}\equiv2^{\,s-3}$ etc.).
\end{lemma}

\begin{proof}
$k_0<q$ gives $s\le 7$; $qs$ odd gives $s$ odd. Then
$4k_0=(qs-1)/2=s\cdot\frac{q-1}{2}+\frac{s-1}{2}$, and
$2^{(q-1)/2}\equiv(\frac{2}{q})\pmod q$ is Euler's criterion; the
Legendre symbol is read off $q\bmod 8$.
\end{proof}

\begin{lemma}[Fermat chain constants]\label{lem:fermat}
For an odd prime $q$, $\;2^{\,4q}\equiv 2^{4}=16\pmod q$, and likewise
$2^{-4q}\equiv 16^{-1}\pmod q$.
\end{lemma}

\begin{theorem}[sieve-batched exponentiation]\label{thm:sieve}
All $2$-power residues with $k$-linear exponents, over all terms of
\eqref{eq:normal} with $k\le K$, are computable in
$O(K\log\log K)$ word operations; all fixed-exponent residues
($5^{\,n}\bmod\mu$, $2^{\,n+O(1)}\bmod\mu$) likewise in
$O(K\log\log K)$ word operations, the per-term cost of the
once-per-prime-power computations vanishing as $O(1/\log K)$. Hence the
entire exponentiation phase costs $O(N\log\log N)$ word operations.
\end{theorem}

\begin{proof}[Proof sketch]
A segmented sieve over each progression delivers the factorization of
every $\mu$ at $O(K\log\log K)$ total cost. Fix a prime power $P=q^{e}$
dividing some moduli; the indices $k$ with $P\mid \mu(k)$ form an
arithmetic progression mod $P$, along which $2^{\,4k}\bmod P$ is a
geometric sequence with ratio $2^{\,4P}\bmod P$: by
Lemmas~\ref{lem:euler}--\ref{lem:fermat} the starter and (for $e=1$) the
ratio are $O(1)$, and each further element costs one multiplication. The
number of (term, prime-power) incidences is
$\sum_{k\le K}\Omega'(\mu(k))=O(K\log\log K)$ by Mertens; prime powers
with $e\ge2$ have $O(\sum_q K/q^{2})=O(K)$ incidences and afford honest
exponentiations for their rare starters. Fixed-exponent residues are
computed once per \emph{distinct} prime power --- there are
$O(K/\log K)$ of them, at $O(\log N)$ each --- and every term's residue is
assembled from its $O(\log\log K)$ cached prime-power components by CRT
at $O(1)$ multiplications per incidence.
\end{proof}

\begin{remark}[application to binary BBP]
The classical BBP loop computes $16^{\,d-k}\bmod(8k+j)$ per term by
square-and-multiply, and has since 1997. The structure required by
Theorem~\ref{thm:sieve} --- fixed base, exponent linear in $k$, modulus on
an arithmetic progression in $k$ --- is exactly the structure of that
loop, so the same sieve-batching applies to hexadecimal/binary digit
extraction of $\pi$ and to every BBP-type computation. We are not aware
of a prior appearance of this amortization. Our validation run
(Section~\ref{sec:impl}) measures $2.7\times$ fewer multiplications
already at $K=2\cdot10^{4}$; the ratio scales as
$\Theta(\log K/\log\log K)$.
\end{remark}

\subsection{The B\'ezout split and the $2$-adic normal form}\label{sec:bezout}

\begin{lemma}[B\'ezout split]\label{lem:bezout}
Let $\mu$ be odd, $J\ge1$, $u=2^{-J}\bmod\mu$, and
$w=\mu^{-1}\bmod 2^{J}$. Then, exactly,
\[
\Bigl\{\frac{X}{\mu 2^{J}}\Bigr\}
=\Bigl\{\frac{(X\bmod\mu)\,u\bmod\mu}{\mu}
+\frac{Xw\bmod 2^{J}}{2^{J}}\Bigr\}.
\]
\end{lemma}

\begin{proof}
With $v=(1-u2^{J})/\mu\in\Z$ we have $u2^{J}+v\mu=1$, hence
$X/(\mu2^{J})=Xu/\mu+Xv/2^{J}$; reduce each part mod $1$ and note
$v\equiv\mu^{-1}\pmod{2^{J}}$.
\end{proof}

The first summand is word arithmetic, and its ingredients
$X\bmod\mu$ and $u=2^{-J}\bmod\mu$ are exactly what
Theorem~\ref{thm:sieve} delivers at amortized $O(\log\log N)$: \emph{the
odd halves of all band terms join the sieved word phase.} The dyadic
halves collapse:

\begin{lemma}[dyadic consolidation]\label{lem:consolidate}
Let $J_{\max}=\max_t J_t$ and
\[
V \;=\; \sum_t \pm\, 2^{\,J_{\max}-J_t}\,
\bigl(\mu_t^{-1}\bmod 2^{\,J_{\max}}\bigr) \;\bmod\; 2^{\,J_{\max}} .
\]
Then
$\sum_t \pm\bigl\{X(\mu_t^{-1}\bmod 2^{J_t})/2^{J_t}\bigr\}
\equiv\bigl\{XV/2^{\,J_{\max}}\bigr\}\pmod 1$.
\end{lemma}

\begin{proof}
Write $\mu^{-1}\bmod 2^{J_{\max}}=(\mu^{-1}\bmod 2^{J})+2^{J}t$; then
$X(\mu^{-1}\bmod 2^{J_{\max}})/2^{J}=X(\mu^{-1}\bmod 2^{J})/2^{J}+Xt$
and the integer $Xt$ vanishes mod $1$. Summing and factoring
$2^{-J_{\max}}$ gives the claim; replacing $V$ by $V\bmod 2^{J_{\max}}$
changes the value by an integer multiple of $X$.
\end{proof}

\paragraph{The refined algorithm.}
The extractor is thereby restructured as:
\begin{enumerate}
\item \textbf{Sieve phase} (Theorem~\ref{thm:sieve}): factorizations,
chains, caches --- $O(N\log\log N)$ word operations.
\item \textbf{Word phase}: all easy terms and all odd halves of band
terms accumulate in fixed point --- $O(N\log\log N)$ word operations.
\item \textbf{$V$ phase}: build $V$ by a balanced merge of
$(\text{numerator},\text{denominator-product})$ pairs \emph{modulo powers
of two}: node arithmetic is truncated multiplication only, and one final
Newton--Hensel inversion mod $2^{J_{\max}}$ recovers $V$ --- cost
$O(\M(N\log N)\log N)$, division-free. (The exact-sum form
$X\cdot(A/B)$ with $A/B=\sum_t\pm1/(\mu_t2^{J_t})$ built by binary
splitting is an equivalent second variant with the same asymptotics.)
\item \textbf{Finish}: one product $XV\bmod 2^{J_{\max}}$, one window
read, digit output --- $O(\M(N))$.
\end{enumerate}
Every logarithm in Theorem~\ref{thm:main} now lives in step 3.

\section{Generalizations}\label{sec:general}

\begin{theorem}\label{thm:general}
Let $\xi$ admit a BBP-type formula in base $2^{r}$,
$\xi=\sum_{k\ge0}2^{-rk}\,p(k)/q(k)$ with integer polynomials and $q$ having
word-size values, in the sense of~\cite{BBP97}. Then the $N$-th decimal
digit of $\xi$ is computable in the time and space of
Theorem~\ref{thm:main}.
\end{theorem}

\begin{proof}
Termwise, $10^{\,n}2^{-rk}p(k)/q(k)=2^{\,n-rk}5^{\,n}p(k)/q(k)$; folding the
$2$-part of $q(k)$ into the exponent yields the normal form
$\pm2^{E}5^{\,n}p(k)/\mu$ with $\mu$ odd, and Lemmas
\ref{lem:band}--\ref{lem:h} apply with numerator $p(k)X$ ($h$ and $u$ scale
by the word-size $p(k)$ before the fixed-point division). All counts are as
before.
\end{proof}

For example $\ln 2=\sum_{k\ge1}1/(k2^{k})$ gives, with $k=2^{s}\mu$,
terms $2^{\,n-k-s}5^{\,n}/\mu$: the first quasi-linear-time decimal digit
extractor for $\ln 2$. The same holds for $\pi^{2}$, $\ln^{2}2$, Catalan-type
combinations, and the rest of the base-$2^r$ BBP zoo of~\cite{BaileyComp}.

\begin{corollary}[any base]\label{cor:anybase}
For any base $B\ge 2$, the $N$-th base-$B$ digit of any base-$2^{r}$ BBP
constant is computable in $\Ot(N)$ time with $\Theta(N\log B)$ bits of
read-only auxiliary space, namely the expansion of $Q^{\,N-1}$ where $Q$ is
the odd part of $B$.
\end{corollary}

\begin{proof}
$B^{\,n}=2^{\,an}Q^{\,n}$ with $Q$ odd; the normal form becomes
$\pm2^{E}Q^{\,n}/\mu$ and the argument is unchanged with $X=Q^{\,n}$.
(For odd $B$, $a=0$ and every non-easy term is band or tail.)
\end{proof}

In particular the base-$5$ digits of $\pi$ are computable in quasi-linear
time with a shared string $X=5^{\,N-1}$ --- in that case \emph{every} term
with $k\ge1$ is a band or tail term. This answers the question left open
in~\cite{Zudilin24} affirmatively at $\Theta(N)$-bit space: the terms on
which his identity~(3) fails (denominator $5^{2n+1-d+k}m$) are exactly our
band terms, and Lemma~\ref{lem:band} computes them. Whether polylogarithmic
space suffices remains open, and our reliance on stored middle bits of
$5^{\,n}$ suggests a connection to the apparent hardness of computing
isolated middle bits of integer powers.

We also note that the framework accepts any Machin-type formula whose
arctangent arguments are $2$-$5$-smooth, such as the 1844
Strassnitzky--Dase identity
$\pi/4=\arctan\frac12+\arctan\frac15+\arctan\frac18$
(equivalently $(2+i)(5+i)(8+i)=65(1+i)$ in $\mathbb{Z}[i]$); by
St{\o}rmer-type finiteness these smooth identities are scarce, and we use
this one only as an independent correctness check.

\section{The barrier}\label{sec:barrier}

After Section~\ref{sec:improve}, every phase except the construction of
$V$ runs in $O(N\log\log N)$ word operations; the $V$ phase remains at the
product-tree bound. This section characterizes the obstruction.

\subsection{Status of lower bounds}

No unconditional lower bound is available at any level: proving even
$\Omega(N)$ for the digit itself would require ruling out closed forms for
the digits of $\pi$, far beyond current techniques --- an obstacle shared
with binary BBP. Within the framework, $\Omega(N)$ is forced structurally:
$\Theta(N)$ terms each contribute at $O(1)$ magnitude modulo $1$, and the
$u$-windows jointly tile $\Theta(N)$ bits of $X$. The natural conjectural
floor is $N\log N$: every known route to the string of $X=5^{\,N-1}$ costs
$\Theta(\M(N))$, and $\M(N)=\Theta(N\log N)$ is believed optimal --- the
upper bound is Harvey--van der Hoeven~\cite{HvdH21}, and a matching lower
bound is known only conditionally, on a network-coding
conjecture~\cite{AFKL19}; no reduction from multiplication to fixed-base
powering is known, so even that conditional bound does not formally
transfer. Sublinear time would require abandoning the series paradigm
altogether, since every term matters at full weight modulo $1$ --- and
three decades of work on binary BBP have not produced $o(N)$ extraction in
any base. Everything between the proven $O(N\log^{3}N)$ and the
conjectural $N\log N$ floor is the subject of what follows.

\subsection{Equivalent faces}

The denominator product entering the $V$ phase is, per family, a
rising-factorial-type product: e.g.
$\prod_{k\le K}(8k+1)=8^{\,K+1}(1/8)^{\overline{K+1}}$. Computing such
products (or $V$, which every known merge scheme routes through them)
modulo $2^{\Theta(N)}$ below the product-tree bound
$\M(n\log n)\log n$ is precisely the classical factorial barrier: no
algorithm is known to compute $n!$ faster, despite long
attention~\cite{Borwein85,Strassen77,BGS07}.

A second face: $(\Z/2^{m}\Z)^{*}=\langle-1\rangle\times\langle3\rangle$,
so $\prod_{k\le K}(1+8k)\equiv\pm 3^{E}\pmod{2^{m}}$ where
\[
E=\sum_{k\le K}\operatorname{dlog}_3(1+8k)\ \bmod\ 2^{\,m-2},
\qquad m=\Theta(N).
\]
Computing the \emph{single integer} $E$ is thus a distillation of the
multiplicative side of the barrier --- and it falls
(Theorem~\ref{thm:crack} below). Expanding
$\operatorname{dlog}_3(1+8k)=\log_2(1+8k)/\log_2 3$ through the $2$-adic
logarithm turns $E$ into $\sum_j c_j S_j(K)$ with power sums
$S_j(K)=\sum_{k\le K}k^{j}$; Faulhaber's formula then prices the exchange
in Bernoulli numbers, and the coefficient mass of every arrangement we
tried is $\Theta(N^{2})$ bits. Third face: merge depth --- summing $T$
inhomogeneous word-size rationals below $\M(\text{total})\log T$.

A fourth face is analytic, and rests on the observation that
Euler--Maclaurin summation, divergent-asymptotic over $\mathbb{R}$, is
\emph{exactly convergent} over $\Z_2$ for our summand.

\begin{lemma}[$2$-adic Euler--Maclaurin]\label{lem:em}
Let $f(x)=\operatorname{Log}_2(1+8x)$ (the $2$-adic logarithm) and
$F(x)=\bigl((1+8x)(f(x)-1)+1\bigr)/8$. Then, as an identity in $\Z_2$
with the right side absolutely convergent (the $i$-th correction has
$2$-adic valuation $\ge 6i-8$),
\[
\sum_{k=1}^{K}f(k)\;=\;F(K)+\tfrac{1}{2}f(K)
+\sum_{i\ge1}\frac{B_{2i}}{(2i)!}
\Bigl(f^{(2i-1)}(K)-f^{(2i-1)}(0)\Bigr),
\]
where $f^{(r)}(x)=(-1)^{r-1}(r-1)!\,8^{r}(1+8x)^{-r}$.
(Machine-validated exactly at $(K,m)=(100,160)$ and $(257,224)$ $2$-adic
bits.)
\end{lemma}

Since the correction series is, up to elementary factors, the Stirling
tail $\sum_i B_{2i}\,w^{2i-1}/\bigl(2i(2i-1)\bigr)$ evaluated at
$w=8/(1+8K)$ (valuation $3$), the lemma says: $E$ --- hence
$\prod_{k\le K}(1+8k)\bmod 2^{m}$, hence the barrier --- \emph{is} a
single-point evaluation of Diamond's $2$-adic log-gamma
function~\cite{Diamond77} in its domain of convergent asymptotics, plus
elementary terms; and conversely, such evaluations at arguments
$8/(1+8K)$, $K=\Theta(m)$, reduce back to $E$-type sums by reading the
identity in reverse. Choosing $K=2^{s}$ and running the dyadic doubling
recursion instead reproduces partial $2$-adic Hurwitz zeta values at
positive integers, confirming the same equivalence from the Iwasawa side
(cf.~\cite{Washington}). Every known evaluation strategy for these
$p$-adic special functions pays the tapered-Bernoulli mass
$\Theta(N^{2})$ bits or falls back to the product tree --- but the
reduction imports the problem into computational Iwasawa theory, where
evaluation of $p$-adic $L$-functions is an active algorithmic topic, and
progress there now transfers here.

A fifth face makes the \emph{input} side provably cheap. The counting
measure $\mu$ on our progression restricted to $[0,2^{t})$ has Amice
transform
\[
A_{\mu}(T)\;=\!\!\sum_{\substack{u<2^{t}\\ u\equiv1\,(8)}}\!\!(1+T)^{u}
\;=\;(1+T)\,\frac{(1+T)^{2^{t}}-1}{(1+T)^{8}-1},
\]
sparse and rational --- $t$ squarings --- and $E$ is the \emph{first
Taylor coefficient} of the Leopoldt $\Gamma$-transform (the pushforward
of $\mu$ along $\operatorname{dlog}_{9}$) of this nearly-free object;
both statements are machine-verified. The $\Gamma$-transform is the
central computation of Iwasawa theory, and every known way of performing
it pays the quadratic mass. Moreover, tracing the state of the
dyadic-doubling computation shows that at every scale the true state is
the node set --- $O(N\log N)$ bits --- while the moment, Bernoulli,
inverse-power-sum and Iwasawa-coefficient representations that carry the
$\Theta(N^{2})$ mass are redundant expansions of it: \emph{no entropy or
output-size argument can defend the barrier at any intermediate point}.
The mass is a property of every known basis, not of the problem. We note
the archimedean contrast: the corresponding real object --- a Stirling
tail at one point --- is computable in quasi-linear time via the Binet
integral and double-exponential quadrature, and the $2$-adic twin resists
precisely for lack of a $p$-adic fast quadrature. This lack is
structural, not an accident of ignorance: translation-invariance makes
the Volkenborn/Riemann-sum scheme the \emph{canonical} $2$-adic
quadrature --- its Euler--Maclaurin acceleration \emph{is} the Bernoulli
mass, Bernoulli numbers being themselves the Volkenborn integrals of
monomials --- whereas over $\mathbb{R}$ the freedom to deform contours
and measures is exactly what powers fast quadrature. The only $2$-adic
deformation freedom is the choice of measure, which is the
Iwasawa-theoretic route already accounted above.

Two further faces close the circle. Sixth: the $\psi$-operator of
Coleman theory intertwines with our dyadic doubling (the measure is
Frobenius-structured, $A_{\mu}$ being a ratio of iterates of
$\varphi(T)=T^{2}+2T$), and exhibits $\mu$ as a level-$t$ truncation of
the Mazur measure whose $\Gamma$-transform is the Kubota--Leopoldt
$2$-adic zeta function: computing $E$ fast is thereby \emph{the same
problem} as computing $\zeta_{2}$ to high $2$-adic precision, an
independently open problem of computational Iwasawa theory. Seventh, and
most elementary: since
$\operatorname{Log}_2 u=\lim_r (u^{2^{r}}-1)/2^{r}$, one has the
machine-validated congruence
\[
\sum_{k\le K}(1+8k)^{2^{r}}\;\equiv\; K \;+\;
2^{r}\!\sum_{k\le K}\operatorname{Log}_2(1+8k)
\pmod{2^{\,2r+5}},
\]
so the entire barrier is equivalent to a statement requiring no $p$-adic
language at all: \emph{compute
$\sum_{k\le K}(1+8k)^{2^{r}}\bmod 2^{2r}$, with $K,r=\Theta(N)$, in
$\Ot(N)$ time.}

\begin{remark}[the analytic root, and why the exits are closed]
\label{rem:root}
The recurring Bernoulli mass has a classical cause: the generating
function $x/(e^{x}-1)$ has infinitely many poles ($2\pi ik$, the poles of
$\zeta$), hence is \emph{not holonomic}, so the Bernoulli sequence
satisfies no polynomial recurrence --- and every fast point-evaluation
technology for series (binary splitting of $P$-recurrences, bit-burst
--- including its $2$-adic version, which achieves quasi-linear precision
for exactly the holonomic class~\cite{CMTV21}) requires such a
recurrence; moreover no bit-burst analogue for the full differentially
algebraic class can exist, since every polynomial-time-computable real is
a value of a polynomial-ODE solution~\cite{BGP17} and the time hierarchy
theorem forbids universal quasi-linear evaluation. Over $\mathbb{R}$ this is harmless
because integral representations route around the coefficients; over
$\Z_2$ the quadrature rigidity above closes that road. The wall is thus
the confluence of two facts: \emph{no recurrence, and no integral
end-run}. Nor can the formula be shopped away: a series whose
denominators were purely geometric would sum to a rational number, so
transcendence forces the polynomial factor in every BBP-type
denominator, which forces moduli sweeping arithmetic progressions, which
forces the faces above. Two further resources are accounted for: quantum
period-finding offers nothing, since the per-element task (one $2$-adic
discrete logarithm) is already classically easy and exact summation to
$2^{-\Theta(N)}$ defeats amplitude-estimation; and the mass is
depth-polylogarithmic parallelizable, so at scale it is a cost in work,
not in wall-clock time.
\end{remark}

\subsection{A crack: the multiplicative faces fall}

\begin{theorem}[sub-product-tree evaluation of the multiplicative faces]
\label{thm:crack}
$\prod_{k\le K}(8k+1)\bmod 2^{m}$ --- and hence $E$, the power sums
$\sum_{k\le K}(1+8k)^{2^{r}}\bmod 2^{2r}$, and the log-gamma values of
Lemma~\ref{lem:em} --- are computable in
$O(\M(N)\log N)=O(N\log^{2}N)$ bit operations.
\end{theorem}

\begin{proof}[Proof sketch]
Borwein's factorial method~\cite{Borwein85} transfers. The exponent of an
odd prime $p$ in the product is
$e_p=\sum_j \#\{k\le K:\,p^{j}\mid 8k+1\}$, a Legendre-type count on an
arithmetic progression, $O(1)$ per prime power after a sieve of cost
$O(N\log\log N)$. The distinct-prime mass is
$\sum_{p\le 8K}\lg p=\Theta(N)$ bits: the factored representation
compresses the $\Theta(N\log N)$-bit term list to $\Theta(N)$ bits, with
multiplicities riding as exponents. Assemble $\prod p^{e_p}\bmod 2^{m}$
by exponent-bit scheduling: per bit level, one truncated product tree
over the participating primes (level masses decay geometrically from
$\Theta(N)$; total $O(\M(N)\log N)$) and one squaring mod $2^{m}$
($O(\log N)$ squarings in all). Finally
$E=\operatorname{Log}_2(\prod)/\operatorname{Log}_2 9$, the $2$-adic
logarithm being quasi-linear in precision~\cite{CMTV21}.
Machine-validated at six sizes through $K=3\cdot10^{6}$: exact agreement
with the direct product, with the measured time ratio growing
monotonically in the factored method's favor --- the log-factor
signature.
\end{proof}

Three consequences. First, the seven faces are \emph{not} constructively
equivalent: they split into a multiplicative cluster --- products, $E$,
power sums, log-gamma/truncated Kubota--Leopoldt values, now all
$O(N\log^{2}N)$, which in particular is the first recorded
sub-product-tree evaluation of Diamond's log-gamma at such arguments (no
$p$-adic Gamma routine exists in the standard libraries at all) --- and
the \emph{additive core}: $V$ itself, the consolidated sum
$\sum_t\pm2^{\,J_{\max}-J_t}\mu_t^{-1}$, whose numerator admits no prime
factorization. Second, Theorem~\ref{thm:main} is unchanged, since the
digit pipeline consumes $V$: the barrier now stands, sharper, on the
additive face alone. Third, a methodological vindication:
Conjecture~\ref{conj:mass} (below) predicted that any crack must leave
every basis it names --- and this one does, entering through the one
representation absent from that list, multiplicities-as-exponents.

\subsection{Routes examined}

\begin{center}
\footnotesize
\setlength{\tabcolsep}{4pt}
\begin{tabular}{ll}
\toprule
Route & Outcome\\
\midrule
lcm-thinned splitting tree & constants only ($\sim$5$\times$); mid-tree mass survives\\
precision-tapered $2$-adic trees & reconverge to the product tree\\
Montgomery batch inversion & quadratic at full precision\\
$2$-adic Lerch series / first-order ODE (bit-burst) & holonomic evaluation, same class\\
Faulhaber--Bernoulli exchange for $E$ & $\Theta(N^{2})$-bit coefficient mass\\
$p$-adic $\Gamma$ at large argument & functional equation unrolls the product\\
Strassen/BGS baby-step--giant-step mod $2^{N}$ & ring-element mass $\Theta(N^{2})$\\
Iwasawa dyadic doubling with inverse-power corrections & initial vector mass $\Theta(N^{2})$\\
$2$-adic Euler--Maclaurin (Lemma~\ref{lem:em}) & exact reduction to $G_{2}$-evaluation; mass persists\\
Newton-series/Abel transform of the $V$-sum & returns hypergeometric splitting\\
$\Gamma$-transform of the sparse-rational counting measure & input free; the pushforward carries the mass\\
dyadic doubling as M\"obius/Taylor shifts & $\Ot(N^{2})$; exposes the operator structure\\
near-linear series composition~\cite{KL24}, checked per face & composition cannot re-index exponents\\
TC$^{0}$/CRT iterated-product circuits~\cite{HAB02} & mod-$p$ dlog mirror: full periods free, partials stall\\
sieve/von Mangoldt restructurings of $E$ & exact rearrangements; the hard core is conserved\\
\bottomrule
\end{tabular}
\end{center}

The Strassen/BGS entry deserves emphasis: $n!\bmod p$ \emph{is}
computable in $\Ot(\sqrt n\,)$ ring operations~\cite{Strassen77,BGS07}
--- sublinear --- because the ring elements are single words. At
precision $2^{\Theta(N)}$ the same technique carries $\Theta(N)$-bit
elements and the advantage vanishes. The barrier is a statement about
\emph{mass}, not about term count.

\subsection{The barrier is not cryptographically protected}

Fast factorials modulo an \emph{arbitrary} integer $m$ would factor $m$
(Shamir's classical argument~\cite{Shamir79}: binary-search
$\gcd(k!\bmod m,\,m)$), which is often cited as evidence that factorials
must be hard. That shield does not cover our instance: the modulus is the
fixed prime power $2^{\Theta(N)}$, where the gcd argument is vacuous. To
our knowledge nothing but the absence of an idea protects
$\prod(8k+1)\bmod 2^{\Theta(N)}$, $E$, or $V$. We consider this the most
attackable formulation and state it as the target:

\begin{openproblem}\label{op:E}
Compute the additive core
$V=\sum_t \pm\,2^{\,J_{\max}-J_t}\bigl(\mu_t^{-1}\bmod 2^{J_{\max}}\bigr)
\bmod 2^{J_{\max}}$ --- equivalently its numerator
$\sum_t \pm\,2^{\,J_{\max}-J_t}\prod_{s\ne t}\mu_s \bmod 2^{J_{\max}}$
--- in $O(\M(N)\log N)$ bit operations. Via the pipeline of
Section~\ref{sec:improve} this immediately yields decimal (and binary)
digit extraction of $\pi$ in $O(N\log^{2}N)$. The denominator product is
computable within that budget (Theorem~\ref{thm:crack}); the \emph{sum}
is the open half, and no factoring-hardness or information-theoretic
argument protects it.
\end{openproblem}

\begin{conjecture}[mass conservation, additive form]\label{conj:mass}
Every algorithm computing the additive core $V$ of
Open Problem~\ref{op:E} by ring operations over $\Z/2^{O(N)}$ whose
intermediate values are coefficient arrays in the bases occurring in
this section --- moments, Bernoulli/Iwasawa coefficients, inverse-power
sums, symmetric functions, or per-prime residue systems --- reads or
writes $\Omega(N^{2}/\log^{O(1)}N)$ bits in total. (The original form of
this conjecture also covered $E$; Theorem~\ref{thm:crack} resolved that
case exactly as the conjecture's guidance predicted --- by an
unlisted representation.)
\end{conjecture}

The routes of the table, spanning real-analytic, $2$-adic, mod-$p$ and
circuit-theoretic bases, are the conjecture's evidence; the compactness
observation above shows it cannot be proven information-theoretically,
and an unconditional proof would amount to an explicit superlinear
circuit lower bound, beyond current techniques. Its practical content is
guidance: a resolution of Open Problem~\ref{op:E} must operate outside
every basis the conjecture names.

\section{Implementation and verification}\label{sec:impl}

Two v1 implementations accompany this paper. The pure-Python reference
\path{decimal_bbp_extract.py} covers both $\pi$ and $\ln2$ and includes
an exact per-term auditor and self-tests. The dependency-free C++ port
\path{three_phase_extract.cpp} adds OpenMP, Barrett-reduced Horner over
the shorter of the prefix $[0,J)$ or suffix $[J,L)$ per band term, a
$96$-bit fixed-point accumulator, and internal computation of $X$ by
word-multiply passes. The v1 loops realize the naive band reduction, hence
carry the $O(N^{2}/w)$ word-operation term. The batching of
Theorem~\ref{thm:batch} is additionally implemented as a Python reference
(\path{extract_v2}: one accumulating remainder tree over the merged band
terms of all four families, with variable chunk widths $\delta_t$, which
also covers $\ln 2$'s non-uniform exponents after sorting by $J$); it
reproduces the v1 fixed-point accumulator bit-for-bit at $88$
position/constant combinations. The refined pipeline of
Section~\ref{sec:improve} itself is implemented as a third path
(\path{extract_v3}: the B\'ezout word phase plus the division-free $V$
merge, with symbolic shifts and per-node precision caps, over GMP
integers via \path{gmpy2} when available); it reproduces the v1 digits
across the same battery and at every benchmark position below.

\paragraph{Correctness.}
All of the following were verified against digits computed independently by
integer Chudnovsky binary splitting (and, for $\ln 2$, by binary splitting
of $2\,\mathrm{artanh}(1/3)$):
positions $1$--$60$ of $\pi$ continuously (8-digit windows);
$\pi$ at $N=10^{2},10^{3},5\cdot10^{3},10^{4},10^{5},10^{6}$;
$\ln2$ at $1$--$40$ and $N=10^{2},10^{3},10^{4},10^{5}$;
agreement at every tested position between the BBP-16 path and the
independent Strassnitzky--Dase path;
and an exact-rational audit of $397$ random band/tail terms (worst
fixed-point deviation $2^{-153.6}$, bound $2^{-142}$).
At $N=10^{6}$ the extractor returns
$130927562832$ at positions $10^{6}..10^{6}{+}11$, matching the reference.
At $N=10^{7}$ the C++ implementation returns $725915133612$ at positions
$10^{7}..10^{7}{+}11$; we verified this window against two independently
published digit corpora --- Google's $100$-trillion-digit computation
(served at \path{pi.delivery}, with the index convention confirmed by
probing the leading digits) and the $200$-million-digit corpus at
\path{angio.net}, whose substring search locates the window at exactly
position $10^{7}$ after the decimal point.

\paragraph{Identity-level validations.}
The mechanisms of Sections \ref{sec:improve}--\ref{sec:barrier} are backed
by exact machine checks: the chain/CRT mechanism of
Theorem~\ref{thm:sieve} reproduces $2^{\,4k}\bmod(8k+1)$ for all
$k\le 2\cdot10^{4}$ (with the Euler-criterion consequence
$2^{4k}\equiv1$ for prime $8k+1$ confirmed throughout, and $2.7\times$
fewer word multiplications than per-term square-and-multiply at that
size); Lemma~\ref{lem:bezout} was verified exactly on $300$ random band
terms; Lemma~\ref{lem:consolidate} was verified exactly against the
termwise dyadic sum; Lemma~\ref{lem:em} was verified exactly at
$(K,m)=(100,160)$ and $(257,224)$; the fifth-face Amice form and the
extraction $E=\operatorname{Log}(\prod u)/\operatorname{Log}9$ were
verified exactly at $96$ bits; and the seventh-face congruence was
verified exactly at $(K,r)=(200,120)$ --- all $245$ bits --- including
recovery of the Log-sum from the elementary power sum alone.

\paragraph{Timings} (Windows 11 desktop, 16 hardware threads; C++ with
OpenMP; Python single-thread, CPython 3.10):

\begin{center}
\begin{tabular}{rrrl}
\toprule
$N$ & Python v1 & C++ v1 & auxiliary string\\
\midrule
$10^{4}$ & $0.12$ s & $0.006$ s & $2.9$ KB\\
$10^{5}$ & $6.4$ s & $0.33$ s & $29$ KB\\
$10^{6}$ & $571$ s & $32.3$ s & $290$ KB\\
$10^{7}$ & --- & $3421$ s & $2.9$ MB\\
\bottomrule
\end{tabular}
\end{center}

For v1 the measured exponent bends from $N^{1.67}$ at small sizes to
$N^{2.03}$ over the decade $10^{6}$--$10^{7}$, as predicted for the naive
band reduction. For an external baseline we compiled the
verification-grade implementation of Gourdon's method --- the strongest
published competitor --- and raced it on the same machine: it measures
$0.59$\,s, $32.3$\,s and $2069$\,s at $N=10^{4},10^{5},10^{6}$ (its
digits agreeing with our verified windows throughout), so the refined
pipeline below is respectively $10\times$, $42\times$ and $232\times$
faster, with the gap growing without bound. The refined pipeline settles
the practical question:

\begin{center}
\begin{tabular}{rrrr}
\toprule
$N$ & v1 Python & v1 C++ (16 threads) & v3 Python (1 thread, GMP)\\
\midrule
$10^{4}$ & $0.12$ s & $0.006$ s & $0.06$ s\\
$10^{5}$ & $6.4$ s & $0.33$ s & $0.77$ s\\
$10^{6}$ & $571$ s & $32.3$ s & $8.9$ s\\
$10^{7}$ & --- & $3421$ s & $111.6$ s\\
\bottomrule
\end{tabular}
\end{center}

The v3 exponent fitted over $10^{4}$--$10^{7}$ --- three decades --- is
$N^{1.077}$, with no upward bend; the quasi-linear regime of
Theorem~\ref{thm:main} is thus visible directly, in interpreted code. At
$N=10^{7}$ the single-threaded v3 outperforms the $16$-thread compiled v1
by $30.7\times$ (and returns the corpus-verified window
$725915133612$); against a hypothetical quadratic continuation of v1's
own small-$N$ constants, the gap at $10^{7}$ exceeds two orders of
magnitude. The v3 phase split at $10^{7}$ ($64$ s exponentiation loop,
$47$ s for the entire $V$ phase, $0.05$ s to build $X$) confirms that the
plain word loop is now the dominant cost --- precisely the phase that the
sieve-batching of Theorem~\ref{thm:sieve} addresses. A pure-Python
integration of that batching (v4) validates the mechanism at scale ---
identical digits through $N=10^{7}$, with the chain/CRT machinery
processing $3.3\cdot10^{6}$ terms --- but is $1.35\times$ \emph{slower}
there in wall clock: the amortization's gain is in multiplication count,
which interpreted per-statement overhead conceals ($\approx$8--10
interpreted operations per chain hit against a single C-level
exponentiation call). Realizing Theorem~\ref{thm:sieve} in running time,
as opposed to operation count, therefore awaits the compiled port, for
which v4 serves as the correctness oracle.

\section{Open problems}

(1) Integrate the sieve-batching of Theorem~\ref{thm:sieve} into the v3
pipeline --- its word loop is now the measured bottleneck --- and port to
C+GMP for record-scale positions. (2) Open Problem~\ref{op:E}: the one remaining logarithm,
in any of its four faces. (3) Decimal digit extraction in sublinear
space --- the original problem of~\cite{BBP97} --- remains open; is there
a lower-bound obstruction linking it to the complexity of isolated middle
bits of $5^{\,n}$? (4) The per-term independence invites a distributed
record attempt: verifying isolated \emph{decimal} positions of record
computations of $\pi$ (a practical need noted by record holders) requires
exactly this primitive. (5) Zudilin's Gaussian-integer series can
presumably be treated by the same band identity over $\mathbb{Z}[i]$ with
a shared string for $(1{-}2i)^{\,d}$; carrying this out would give a
second, structurally different quasi-linear base-5 extractor.

\subsection*{Acknowledgments}
The author thanks the maintainers of the computational records and surveys
cited here. % [venue-dependent AI-assistance disclosure may go here]

\begin{thebibliography}{99}

\bibitem{BBP97}
D.~Bailey, P.~Borwein, S.~Plouffe,
\emph{On the rapid computation of various polylogarithmic constants},
Math.\ Comp.\ \textbf{66} (1997), 903--913.

\bibitem{BBG04}
D.~Borwein, J.~M.~Borwein, W.~F.~Galway,
\emph{Finding and excluding $b$-ary Machin-type individual digit formulae},
Canad.\ J.\ Math.\ \textbf{56} (2004), 897--925.

\bibitem{Plouffe96}
S.~Plouffe,
\emph{On the computation of the $n$'th decimal digit of various
transcendental numbers}, 1996; arXiv:0912.0303.

\bibitem{Bellard97}
F.~Bellard,
\emph{Computation of the $n$'th digit of pi in any base in $O(n^{2})$},
1997, \url{https://bellard.org/pi/pi_n2/pi_n2.html}.

\bibitem{Gourdon03}
X.~Gourdon,
\emph{Computation of the $n$-th decimal digit of $\pi$ with low memory},
unpublished manuscript, February 11, 2003,
\url{http://numbers.computation.free.fr/Constants/Algorithms/nthdecimaldigit.pdf}.

\bibitem{Plouffe22}
S.~Plouffe,
\emph{A formula for the $n$'th decimal digit or binary of $\pi$ and powers
of $\pi$}, arXiv:2201.12601 (2022).

\bibitem{Zudilin24}
W.~Zudilin,
\emph{A BBP-style computation for $\pi$ in base 5},
arXiv:2409.10097 (2024). (v1 titled ``\ldots in base 10''.)

\bibitem{BaileyComp}
D.~H.~Bailey,
\emph{A compendium of BBP-type formulas for mathematical constants},
manuscript, updated 2023,
\url{https://www.davidhbailey.com/dhbpapers/bbp-formulas.pdf}.

\bibitem{BorodinMoenck74}
A.~Borodin, R.~Moenck,
\emph{Fast modular transforms},
J.~Comput.\ System Sci.\ \textbf{8} (1974), 366--386.

\bibitem{Bernstein04}
D.~J.~Bernstein,
\emph{Scaled remainder trees}, manuscript, 2004,
\url{https://cr.yp.to/arith/scaledmod-20040820.pdf}.

\bibitem{CGH14}
E.~Costa, R.~Gerbicz, D.~Harvey,
\emph{A search for Wilson primes},
Math.\ Comp.\ \textbf{83} (2014), 3071--3091.

\bibitem{HS14}
D.~Harvey, A.~V.~Sutherland,
\emph{Computing Hasse--Witt matrices of hyperelliptic curves in average
polynomial time},
LMS J.\ Comput.\ Math.\ \textbf{17} (2014), 257--273.

\bibitem{HvdH21}
D.~Harvey, J.~van der Hoeven,
\emph{Integer multiplication in time $O(n\log n)$},
Ann.\ of Math.\ \textbf{193} (2021), 563--617.

\bibitem{AFKL19}
P.~Afshani, C.~B.~Freksen, L.~Kamma, K.~G.~Larsen,
\emph{Lower bounds for multiplication via network coding},
in: ICALP 2019, LIPIcs \textbf{132}, 10:1--10:12.

\bibitem{ZZ15}
D.~Zeilberger, W.~Zudilin,
\emph{The irrationality measure of $\pi$ is at most $7.103205334137\ldots$},
Exp.\ Math.\ \textbf{24} (2015), 419--423.

\bibitem{Borwein85}
P.~B.~Borwein,
\emph{On the complexity of calculating factorials},
J.\ Algorithms \textbf{6} (1985), 376--380.

\bibitem{Strassen77}
V.~Strassen,
\emph{Einige Resultate \"uber Berechnungskomplexit\"at},
Jahresber.\ Deutsch.\ Math.-Verein.\ \textbf{78} (1976/77), 1--8.

\bibitem{BGS07}
A.~Bostan, P.~Gaudry, \'E.~Schost,
\emph{Linear recurrences with polynomial coefficients and application to
integer factorization and Cartier--Manin operator},
SIAM J.\ Comput.\ \textbf{36} (2007), 1777--1806.

\bibitem{Shamir79}
A.~Shamir,
\emph{Factoring numbers in $O(\log n)$ arithmetic steps},
Inform.\ Process.\ Lett.\ \textbf{8} (1979), 28--31.

\bibitem{CMTV21}
X.~Caruso, M.~Mezzarobba, N.~Takayama, T.~Vaccon,
\emph{Fast evaluation of some $p$-adic transcendental functions},
in: ISSAC 2021; arXiv:2106.09315.

\bibitem{BGP17}
O.~Bournez, D.~S.~Gra\c{c}a, A.~Pouly,
\emph{Polynomial time corresponds to solutions of polynomial ordinary
differential equations of polynomial length},
J.~ACM \textbf{64} (2017), no.~6, art.~38.

\bibitem{KL24}
Y.~Kinoshita, B.~Li,
\emph{Power series composition in near-linear time},
in: FOCS 2024.

\bibitem{HAB02}
W.~Hesse, E.~Allender, D.~A.~Mix Barrington,
\emph{Uniform constant-depth threshold circuits for division and
iterated multiplication},
J.~Comput.\ System Sci.\ \textbf{65} (2002), 695--716.

\bibitem{Diamond77}
J.~Diamond,
\emph{The $p$-adic log gamma function and $p$-adic Euler constants},
Trans.\ Amer.\ Math.\ Soc.\ \textbf{233} (1977), 321--337.

\bibitem{Washington}
L.~C.~Washington,
\emph{Introduction to Cyclotomic Fields}, 2nd ed., GTM 83,
Springer, 1997.

\end{thebibliography}

\end{document}
